\begin{textsample}{POS Dim 1 – vem\_ed\_26 – Score 1.00 – vem\_ed\_26/t139.txt}  \label{ex:f1_pos_050}
Visitas internacionais No mês de outubro, duas visitas internacionais ocorreram nas Fatecs Itaquera e São Paulo, na capital, como fruto da realização de PCIs / Cesu entre essas Fatecs e instituições dos EUA e Colômbia.
No dia 16, a Fatec Itaquera recebeu a visita de dois representantes do Kirkwood Community College, instituição de ensino do estado de Iowa, Estados Unidos: Joseph Greathouse, Diretor do Departamento de Tecnologias Industriais, e Dane Dermody, professor assistente de Tecnologias de Soldagem Avançadas.
A visita foi motivada pelo PCI / Cesu realizado com estudantes de Soldagem do Kirkwood Community College e da Fatec Itaquera e teve como propósito um estreitamento de laços entre as instituições para futuros intercâmbios de alunos e professores. “A visita foi muito apreciada e esperamos continuar nossa colaboração no futuro”, comentou Anderson Ribeiro, coordenador do curso de Fabricação Mecânica, em postagem da página da Fatec Itaquera no LinkedIn.
Nos dias 29 e 30, um grupo de 22 alunos e professores da Unisimón (Colômbia) visitou a Fatec São Paulo.
Os preparativos para a visita presencial foram articulados em um PCI / Cesu realizado entre as duas instituições.
A equipe dos PCIs / Cesu auxiliou a estruturar a programação e acompanhou a visita.
Estudantes de Secretariado da Fatec São Paulo, orientados pela professora Glauce Gomes, organizaram a \textbf{recepção} e o coffee break.
O diretor da Fatec São Paulo, Josué Souza de Gois, e a coordenadora do curso de Secretariado, Maria do Carmo Ferreira Lima, deram as boas-vindas aos visitantes.
Douglas Dias, coordenador do curso de Gestão de Empreendimentos Gastronômicos, ministrou uma palestra, no dia 29, sobre a cultura gastronômica em São Paulo.
Em seguida, o grupo fez uma visita técnica e almoçou no Pão do Povo da Rua, projeto premiado pela Organização das Nações Unidas para Alimentação e Agricultura (FAO) em outubro de 2024: http://www.paodopovodaarua.com.br/. continuação O projeto Pão do Povo da Rua distribui 3.000 pães e 1.500 refeições diariamente a moradores de rua em São Paulo, além de oferecer capacitações em panificação, bolsa-auxílio e assistência social para pessoas em situação de rua.
Depois do almoço, alunos e professores da Fatec São Paulo e da Unisimón fizeram uma visita monitorada ao Museu da Língua Portuguesa.
No dia 30, Lucas Martins, professor do curso de Produção Cultural, fez uma palestra sobre roteiro para videogames.
A coordenadora do curso, Joana Ormundo, conduziu uma roda de conversa com os estudantes.
Após o almoço, André Guilles (Cesu), ofereceu uma oficina de criatividade.

% matched lemmas: recepção
\end{textsample}
